\documentclass[11pt,a4paper]{article}
\usepackage[utf8]{inputenc}
\usepackage[T1]{fontenc}
\usepackage[portuguese]{babel}
\usepackage{xcolor}
\usepackage{hyperref}
\usepackage{geometry}
\usepackage{tcolorbox}

\geometry{margin=2.5cm}

\definecolor{primary}{RGB}{39, 174, 96}
\definecolor{secondary}{RGB}{44, 62, 80}
\definecolor{codebg}{RGB}{240, 240, 240}

\hypersetup{colorlinks=true, linkcolor=primary, urlcolor=primary}

\title{\color{secondary}\Huge\textbf{Solução: Exercício 4.3 - Prometheus}}
\author{\color{primary}DevOps com Kubernetes -- 2026}
\date{}

\begin{document}
\maketitle

\section*{\color{primary}Guia de Solução}
\begin{tcolorbox}[colback=green!5!white,colframe=green!60!black,title=Resolução Passo a Passo]
### Passo a Passo

1. \textbf{Expor o Prometheus}: Com a \textit{stack} kube-prometheus instalada via Helm, utilize o \texttt{port-forward} para trazer a porta do servidor localmente: \texttt{kubectl port-forward -n monitoring svc/prometheus-server 9090:80}.
2. \textbf{Acesso e Consulta}: Acesse \texttt{http://localhost:9090} no seu navegador.
3. \textbf{Escrever a Query}: Na barra de buscas do Prometheus, insira a \textit{query} PromQL adequada. Para listar o número de pods criados por StatefulSets no namespace do Prometheus, você pode usar a métrica \texttt{kube_pod_info} filtrada: \texttt{count(kube_pod_info{created_by_kind="StatefulSet", namespace="prometheus"})}. O resultado gráfico deve apresentar o valor 3, correspondendo aos réplicas do StatefulSet em execução.
\end{tcolorbox}

\end{document}
